\chapter{Utilizzo del client testuale e del server}
\label{chap:client-server}

\section{Scopo della modalita' client-server}

La coppia \module{qtServer}/\module{qtClient} rappresenta la modalita' remota del
progetto. A differenza della GUI, che ingloba il backend e lavora in locale, qui
l'algoritmo QT viene eseguito esclusivamente sul server. Il client si limita a stabilire
una connessione, inviare comandi e visualizzare le risposte testuali ricevute. Questa
modalita' e' utile per comprendere l'evoluzione storica del progetto e per utilizzare il
backend senza passare dall'interfaccia grafica.

\section{Avvio del server}

Il server viene avviato tramite la classe \module{server.MultiServer}. Se non viene
specificata alcuna porta, il valore di default e' 8080. All'avvio il server apre una
\module{ServerSocket}, stampa un messaggio di disponibilita' e resta in attesa di nuove
connessioni. Per ogni client viene creato un thread \module{ServerOneClient} indipendente.

\begin{commandblock}
  make run-server PORT=8080
\end{commandblock}

Questa organizzazione implica che piu' client possano collegarsi contemporaneamente senza
condividere lo stesso stato applicativo. Ogni thread mantiene infatti il proprio dataset
corrente e l'ultimo clustering calcolato.

\section{Avvio del client}

Il client testuale richiede esplicitamente due argomenti: indirizzo IP e porta del
server. Una volta avviato, costruisce uno \module{ObjectOutputStream} e uno
\module{ObjectInputStream}, quindi presenta un menu numerico. La connessione resta aperta
finche' l'utente non seleziona l'opzione di uscita.

\begin{commandblock}
  make run-client IP=localhost PORT=8080
\end{commandblock}

All'atto pratico il client mostra un menu con cinque opzioni numerate da 0 a 4. La logica
dell'applicazione dipende fortemente dalla sequenza con cui queste opzioni vengono usate,
perche' alcune richiedono che il server abbia gia' memorizzato un dataset o un clustering
all'interno della stessa sessione.

\section{Descrizione dettagliata dei comandi}

\subsection{Comando 0: caricamento di una tabella da database}

L'opzione 0 chiede il nome della tabella da caricare e invia tale informazione al server.
Il thread \module{ServerOneClient} costruisce quindi un oggetto \module{Data(tableName,
true)}, che usa la connessione JDBC predefinita e recupera le tuple dalla tabella
indicata. In caso di successo il dataset viene memorizzato nello stato della sessione e
qualsiasi clustering precedente viene invalidato.

Dal punto di vista operativo, questo comando e' il prerequisito naturale del comando 1,
che esegue il clustering sul dataset appena caricato. Se la tabella non esiste o il
database non e' raggiungibile, il server restituisce un messaggio di errore testuale
invece della risposta \module{OK}.

\subsection{Comando 1: esecuzione del clustering sulla tabella corrente}

L'opzione 1 richiede il valore di \emph{radius} e lo inoltra al server. Prima di
procedere, il thread verifica che un dataset sia gia' stato caricato nella sessione. Se
manca, il server risponde con un messaggio che invita a usare prima il comando 0.

Quando il dataset e' presente, il server istanzia un nuovo \module{QTMiner}, esegue
\module{compute(data)} e restituisce al client tre elementi: la stringa \module{OK}, il
numero di cluster trovati e una rappresentazione testuale dell'intero
\module{ClusterSet}. Il client visualizza quindi l'output nel terminale. A differenza
della GUI, qui non sono disponibili grafici, dialoghi statistici o funzioni di
esplorazione interattiva del risultato.

\subsection{Comando 2: salvataggio dei cluster su file}

L'opzione 2 usa l'ultimo \module{QTMiner} disponibile nella sessione e salva il
clustering in un file \module{.dmp}. Il nome del file viene costruito automaticamente dal
server concatenando il nome della tabella corrente e un timestamp nel formato
\module{yyyyMMdd\_HHmmss}. Se l'utente prova a usare il comando senza aver prima eseguito
un clustering, il server risponde con un messaggio d'errore esplicito.

Questa scelta automatica del nome file riduce il numero di parametri da inserire nel
client, ma richiede all'utente di prestare attenzione alla cartella di lavoro da cui il
server e' stato avviato, poiche' il file viene scritto proprio in quel contesto.

\subsection{Comando 3: caricamento di un clustering da file}

L'opzione 3 consente di riaprire un clustering gia' serializzato. Il client chiede il
nome del file senza estensione \module{.dmp}, invia tale nome al server e, per ragioni di
compatibilita', chiede anche un radius che sul lato server viene pero' ignorato. Il
\module{QTMiner} viene ricostruito a partire dal file, il dataset associato viene
recuperato e il server restituisce nuovamente una rappresentazione testuale del clustering.

L'uso di questo comando non richiede che un database sia attivo, a patto che il file
\module{.dmp} esista gia' e sia leggibile dal server. Di conseguenza la modalita'
client-server puo' essere usata anche per consultare risultati archiviati in precedenza.

\subsection{Comando 4: chiusura della sessione}

L'opzione 4 non invia un'operazione complessa al server, ma termina il ciclo principale
del client e chiude la connessione. Il server rileva la fine della comunicazione come
\module{EOF} e provvede a liberare le risorse relative al thread dedicato.

\section{Stato di sessione e dipendenze tra operazioni}

Il comportamento del client puo' essere compreso pienamente solo tenendo presente che il
server mantiene uno stato per ogni connessione. Tre campi sono particolarmente
importanti: il dataset corrente, l'ultimo miner calcolato e il nome della tabella
caricata. Questa progettazione ha conseguenze immediate:

\begin{itemize}
  \item il comando 1 dipende dal comando 0, salvo che non si stia lavorando su un file
    \module{.dmp};
  \item il comando 2 dipende dal comando 1, perche' salva l'ultimo clustering calcolato;
  \item il comando 3 puo' reintrodurre sia cluster sia dataset direttamente da disco.
\end{itemize}

In altri termini, il client testuale non e' uno strumento puramente stateless. La singola
sessione costituisce un piccolo spazio di lavoro temporaneo sul server.

\section{Indicazioni pratiche d'uso}

Per un uso corretto della modalita' client-server e' consigliabile seguire una procedura
lineare. Si avvia il server, si apre il client, si carica una tabella dal database, si
esegue il clustering con un \emph{radius} scelto con cura e, solo a valle della verifica
dell'output, si decide se salvare il risultato. Qualora si voglia lavorare esclusivamente
su dati gia' archiviati, e' possibile saltare il caricamento da database e usare
direttamente il comando 3.

Pur essendo semplice, questa interfaccia non guida l'utente in modo ricco come la GUI.
Non ci sono convalide avanzate sul contenuto della tabella, non sono disponibili grafici,
e l'output e' interamente testuale. La sua utilita' principale resta quindi quella di
offrire un canale minimale e trasparente verso il backend remoto.
